\documentclass[11pt,a4paper]{article}
\usepackage[T1]{fontenc}
\usepackage[utf8]{inputenc}
\usepackage[main=italian,provide=*]{babel}
\usepackage{amsmath,amssymb}
\usepackage{geometry}
\geometry{margin=2.5cm}
\usepackage{booktabs}
\usepackage{seqsplit}
\usepackage{hyperref}
\hypersetup{colorlinks=true,linkcolor=blue,citecolor=blue,urlcolor=blue}
\newcommand{\ket}[1]{\lvert #1 \rangle}

\title{Manuale d'uso --- correlazioni dinamiche trimero a catena aperta\\
\large versione a stato esatto (\texttt{prepare\_state})}
\author{Note di lavoro --- Tesi triennale, Samuele Schiavi\\ Relatore: Prof.\ Alessandro Chiesa}
\date{}

\begin{document}
\maketitle

\section*{Scopo}

Questo manuale spiega come eseguire la simulazione delle correlazioni dinamiche
$C_{ij}^{\alpha\beta}(t)$ per il trimero a catena aperta nella sua versione
\emph{originale}: lo stato fondamentale $\ket{\psi_0}$ è preparato sul registro
con le ampiezze esatte (\texttt{QuantumCircuit.prepare\_state}), non con un
ansatz VQE. È la versione da usare quando interessa isolare l'errore di
Trotter e quello statistico (shot), senza mescolarli con un residuo di
fedeltà dell'ansatz variazionale. Per la versione con il circuito VQE reale
(\texttt{W-2qC.K2}), vedi il manuale gemello
\texttt{\seqsplit{manuale\_uso\_correlazioni\_trimero\_catena\_vqe.tex}}.

\section*{Prerequisiti}

\textbf{File necessari nella stessa cartella} (tutti nella fase
\emph{correlazioni dinamiche}, trimero catena):
\begin{itemize}
\item \texttt{trimer\_chain\_exact.py} --- Hamiltoniana \texttt{trimer\_hamiltonian\_dm}.
\item \texttt{\seqsplit{trotter\_trimero\_catena.py}} --- evoluzione $U(t)=e^{-iHt}$ via
Suzuki--Trotter (due bond, opzione \texttt{alternate} disponibile, default \texttt{False}).
\item \texttt{\seqsplit{circuito\_correlazioni\_trimero\_catena.py}} --- il circuito Hadamard test a 4 qubit.
\item \texttt{\seqsplit{validate\_circuito\_correlazioni\_catena.py}} --- riferimento classico esatto (\texttt{scipy.linalg.expm}).
\item \texttt{\seqsplit{validate\_shot\_noise\_trimero\_catena.py}} --- validazione a shot finiti.
\item \texttt{scan81\_trimero\_catena.py} --- scan completo delle 81 combinazioni.
\item \texttt{\seqsplit{generate\_figure\_shotnoise\_catena.py}}, \texttt{\seqsplit{generate\_figure\_scan81\_catena.py}} --- figure.
\end{itemize}

\textbf{Librerie Python}: \texttt{qiskit} (2.x), \texttt{qiskit-aer}, \texttt{numpy},
\texttt{scipy}, \texttt{matplotlib} (solo per le figure).

\begin{quote}
\textbf{Importante, a differenza dell'anello}: questa fase lavora su \textbf{due} punti di
lavoro, non uno solo, tenuti entrambi per confronto perché non ancora unificati con il
relatore (vedi \texttt{\seqsplit{simmetrie\_correlatori\_trimero\_catena.tex}}, nota aperta):
\begin{itemize}
\item \textbf{VQE-DM}: $J=1$, $b=b_c=3.0$, $D=0.15$ --- punto dove è stato identificato
l'ansatz canonico \texttt{W-2qC.K2}.
\item \textbf{S0 (Trotter)}: $J=1$, $b=3.0073414$, $D=0.3$ --- punto provvisorio della fase
Trotter, verificato essere l'argmin del gap a $D=0.3$.
\end{itemize}
Gli script che supportano entrambi i punti li selezionano con un argomento da riga di comando
(\texttt{vqedm} o \texttt{s0}, default \texttt{vqedm} se omesso); altri (i moduli di
validazione) eseguono automaticamente entrambi in sequenza. $t=1.3$, $N=200$ passi di Trotter
in tutta la fase, in entrambi i punti.
\end{quote}

\section*{Passo 1 --- un correlatore singolo, da riga di comando Python}

Per calcolare un correlatore specifico, ad esempio $C_{11}^{yy}(t{=}1.3)$ al punto VQE-DM:
\begin{verbatim}
from circuito_correlazioni_trimero_catena import (
    ground_state, correlator_from_circuit,
)

J, b, D = 1.0, 3.0, 0.15
psi0, E = ground_state(J, b, D)

C = correlator_from_circuit(1, 'y', 1, 'y', t=1.3, N=200,
                             J=J, b=b, D=D, psi0=psi0)
print(C)   # atteso: un numero complesso, |C| dell'ordine di alcuni decimi
\end{verbatim}
Non passare \texttt{ansatz\_params}: di default resta \texttt{None} e la preparazione avviene
con \texttt{prepare\_state(psi0)}, ampiezze esatte. Per il punto S0, usare invece
$b{=}3.0073414$, $D{=}0.3$.

\section*{Passo 2 --- validazione contro il riferimento classico esatto}

\begin{verbatim}
python validate_circuito_correlazioni_catena.py
\end{verbatim}
Esegue \textbf{automaticamente entrambi} i punti di lavoro in sequenza (nessun argomento da
riga di comando). Per ciascuno stampa: confronto su cinque correlatori rappresentativi e tre
valori di $t$ (residuo atteso $\sim10^{-3}$, errore di Trotter puro); zeri predetti a $t=0$
(Cor.\ zero-t0-same/diff, atteso $0$ a precisione macchina); relazione $P_{13}$
($C_{11}^{xx}=C_{33}^{xx}$, entro l'errore di Trotter); convergenza in $N$ (rapporto $\approx2$
raddoppiando $N$).

\section*{Passo 3 --- validazione a shot finiti}

\begin{verbatim}
python validate_shot_noise_trimero_catena.py vqedm
python validate_shot_noise_trimero_catena.py s0
\end{verbatim}
A differenza del Passo 2, questo script richiede l'argomento (\texttt{vqedm} o \texttt{s0}):
esegue un solo punto per lancio, non entrambi. Separa i due contributi di errore (Trotter e
statistico) su $C_{11}^{zz}(t{=}1.3)$: rapporto statistico/Trotter $\sim10\times$ a VQE-DM,
$\sim115\times$ a S0 (vedi
\texttt{\seqsplit{circuito\_correlazioni\_trimero\_catena\_spiegato.tex}}). Salva
\texttt{shot\_noise\_catena\_\{vqedm,S0\}\_results.npz}.

\section*{Passo 4 --- scan completo delle 81 combinazioni}

\begin{verbatim}
python scan81_trimero_catena.py vqedm
python scan81_trimero_catena.py s0
\end{verbatim}
Come il Passo 3, un punto per lancio. Richiede qualche minuto per punto (81 combinazioni
$\times$ 2 parti Re/Im $\times$ (statevector + 8192 shots), con cache di trascompilazione).
Stampa un riepilogo (nessuno zero strutturale confermato per la catena, errori medi/massimi,
correlatori più "ricchi") e salva, per ciascun punto:
\begin{itemize}
\item \texttt{scan81\_catena\_\{vqedm,S0\}\_results.npz} --- matrici $9\times9$ di riferimento
classico, statevector e shot, più le statistiche di errore.
\item \texttt{scan81\_catena\_\{vqedm,S0\}\_results.json} --- stessi dati, riga per riga.
\end{itemize}
\textbf{Attenzione}: percorsi relativi --- va eseguito da dentro la cartella di lavoro.

\section*{Passo 5 --- figure}

\begin{verbatim}
python generate_figure_shotnoise_catena.py
python generate_figure_scan81_catena.py
python generate_circuit_fig_trimero_catena.py
python generate_figures_circuito_trimero_catena.py
\end{verbatim}
I primi due richiedono che i file \texttt{.npz} dei Passi 3--4 (\textbf{entrambi} i punti)
siano già presenti: generano automaticamente una figura per punto. Gli ultimi due generano il
diagramma del circuito e le tre figure di
\texttt{\seqsplit{circuito\_correlazioni\_trimero\_catena\_spiegato.tex}} (convergenza Trotter,
validazione continua, contrasto "nessuno zero sul sito 2") -- questi ultimi lavorano solo al
punto VQE-DM, hard-coded in testa allo script.

\section*{Problemi comuni}

\begin{itemize}
\item \texttt{ModuleNotFoundError} su uno dei moduli locali --- verificare di eseguire dalla
cartella che contiene tutti i file elencati in "Prerequisiti".
\item \texttt{FileNotFoundError} su uno dei file \texttt{scan81\_catena\_*\_results.npz} --- va
prima generato eseguendo \texttt{scan81\_trimero\_catena.py} col punto corrispondente (Passo 4);
non incluso di default perché il calcolo richiede qualche minuto per punto.
\item Risultati leggermente diversi da quelli riportati nei documenti di teoria --- controllare
il punto di lavoro ($J,b,D$): due punti diversi sono in uso in questa fase (vedi riquadro sopra),
a differenza dell'anello che ne ha uno solo. Controllare anche $t,N$: hard-coded in testa a ogni
script.
\item Confusione fra le due topologie --- questo modulo importa da
\texttt{trimer\_chain\_exact.py} e \texttt{trotter\_trimero\_catena.py}, \textbf{non} dagli
omonimi \texttt{\_ring\_}/\texttt{\_anello\_}: le firme delle funzioni differiscono
($J,b,D$ qui, $J,Jp,b,D$ nell'anello --- nessun parametro $Jp$ nella catena).
\end{itemize}

\section*{Riferimenti}

Teoria e derivazioni: \texttt{\seqsplit{simmetrie\_correlatori\_trimero\_catena.tex}},
\texttt{\seqsplit{circuito\_correlazioni\_trimero\_catena\_spiegato.tex}}. Catalogo completo
della fase: \texttt{\seqsplit{trimero\_catena\_correlazioni.tex}}.

\end{document}
